\exercisegroup

\begin{enumerate}
\def\labelenumi{\arabic{enumi})}
\item
  设 P 是 E（R 上的向量空间）中以 0 为顶点的真尖凸锥，p 是 E 上的一个半范数，V 是点 \(x \in E\)（满足 \(p(x) < 1\)）所成之集。设 M 是 E 的一个向量子空间，f 是 M 上的一个线性型。存在 E 上的线性型 g，它延拓 f，在 P 中 \(\geqslant 0\)，且 \(|g(x)| \leqslant p(x)\) 对所有 \(x \in E\) 成立，其充要条件是：对所有 \(x \in M \cap (V + P)\)，有 \(f(x) > -1\)。（为看出条件是充分的，考虑一点 \(x_0 \in M\)（满足 \(f(x_0) = 1\)）、以 0 为顶点由 \(x_0 + V\) 生成的锥 Q，并对赋以这样预序关系的空间 E 应用 prop. 1 的推论（II, 第 21 页）：在该预序下 \(P + Q\) 是诸元素 \(\geqslant 0\) 所成之锥。）
\item
  对一个集 S，设 \(F = \mathcal{B}(S)\) 是 S 中实值有界函数所成的 Banach 空间（I, 第 4 页），并设 M 是赋范空间 E 的一个向量子空间。证明：对 M 到 F 中的每个连续线性映射 f，存在 E 到 F 中的连续线性映射 g，它是 f 的延拓且满足 \(\|g\| = \|f\|\)。
\item
  设 E 是 R 上的向量空间，p 是 E 上的一个次线性函数（II, 第 20 页）。设 A 是凸集，使 \(\inf p(y) > -\infty\)。
\end{enumerate}

\begin{enumerate}
\def\labelenumi{\alph{enumi})}
\tightlist
\item
  证明：函数
\end{enumerate}

\[
q (x) = \inf _ {z \in \mathrm{A}, t \geqslant 0} \left(p (x + t z) - t. \inf _ {y \in \mathrm{A}} p (y)\right)
\]

是 E 上的一个次线性函数，且 \(-p(-x) \leqslant q(x) \leqslant p(x)\)。b) 证明：存在 E 上的线性型 h，使 \(h(x) \leqslant p(x)\) 在 E 中成立，且 \(\inf_{y \in A} p(y) = \inf_{y \in A} h(y)\)（取 h 满足 \(h(x) \leqslant q(x)\)）。

\begin{enumerate}
\def\labelenumi{\arabic{enumi})}
\setcounter{enumi}{3}
\tightlist
\item
  设 A 是 E（R 上的向量空间）的非空子集，p 是 E 上的一个次线性函数。设 B 是由 \(z \in E\) 中使 \(\inf_{x \in A} p(x - z) \leqslant 0\) 的那些 z 所成之集；我们有 \(A \subset B\)，且 \(\inf_{x \in A} p(x) \leqslant p(z)\) 对所有 \(z \in B\) 成立；由此得 \(\inf_{x \in A} p(x) = \inf_{z \in B} p(z)\)。
\end{enumerate}


\begin{enumerate}
\def\labelenumi{\alph{enumi})}
\item
  证明：由 \(y \in \mathbf{E}\) 中使 \(\inf_{z \in \mathbf{B}} p(z - y) \leqslant 0\) 的那些 y 所成之集就是集 \(\mathbf{B}\)。
\item
  由 a) 推出：B 与 E 中任意仿射直线 D 的交在 D 中是闭的（证明：无论 E 的点 a、b 如何，函数 \(t \rightarrow p(a + tb)\) 都在 R 中连续）。
\item
  假设对 A 的每一对点 x、y，存在 \(z \in A\) 使 \(p(z - \frac{1}{2}(x + y)) \leqslant 0\)。证明：对 B 的每一对点 u、v，我们有 \(\frac{1}{2}(u + v) \in \mathbf{B}\)（写出
\end{enumerate}

\[
z - \frac {1}{2} (u + v) = \left(z - \frac {1}{2} (x + y)\right) + \frac {1}{2} (x - u) + \frac {1}{2} (y - v)
\]

其中 \(x, y, z\) 属于 A）。由此推出：B 此时是凸的（利用 \(b\) )）。

\begin{enumerate}
\def\labelenumi{\alph{enumi})}
\setcounter{enumi}{3}
\tightlist
\item
  在 \(c\) ) 的假设下，证明：存在一个线性型 \(h\)（\(\mathbf{E}\) 上的），使 \(h(x) \leqslant p(x)\)，且有 \(\inf_{y \in \mathbf{A}} p(y) = \inf_{y \in \mathbf{A}} h(y)\)（利用 \(c\) ) 与习题 3）。
\end{enumerate}

\begin{enumerate}
\def\labelenumi{\arabic{enumi})}
\setcounter{enumi}{4}
\item
  设 A 是 R 上向量空间 E 的非空子集，p 是 E 上的一个次线性函数。假设对 A 的每一对点 x、y，存在 \(z \in A\) 使 \(p(z - (x + y)) \leqslant 0\)，且 \(p(x) \geqslant 0\) 对所有 \(x \in A\) 成立。证明：存在 E 上的线性型 h，使 \(h(x) \leqslant p(x)\) 在 E 中成立，且 \(h(x) \geqslant 0\) 对 \(x \in A\) 成立。（对诸 \(\frac{1}{n}A\)（n（整数）\(\geqslant 1\)）之并应用习题 4, c)。）
\item
  \begin{enumerate}
  \def\labelenumii{\alph{enumii})}
  \tightlist
  \item
    设 H 是 E（R 上的向量空间）中的一个超平面，p 是 E 上的一个次线性函数。设 f 是 H 上的一个线性型，满足 \(f(y) \leqslant p(y)\) 在 H 中成立。设 a 是 \(\complement H\) 的一点，h 是 E 上延拓 f 且满足 \(h(a) = \inf_{y \in H} (f(y + p(a - y)))\) 的线性型。则 \(h(x) \leqslant p(x)\) 在 E 中成立。证明：对每个线性型 \(g\)（E 上的、延拓 \(f\)、且使 \(g(x) \leqslant p(x)\) 在 E 中成立的），我们也有 \(g(a) \leqslant h(a)\)。
  \end{enumerate}
\end{enumerate}

\begin{enumerate}
\def\labelenumi{\alph{enumi})}
\setcounter{enumi}{1}
\tightlist
\item
  设 V 是 E 的一个向量子空间，f 是 V 上的一个线性型，满足 \(f(y) \leqslant p(y)\) 在 V 中成立。设 S 是 E 的一个非空子集。证明：存在 E 上的线性型 h，它延拓 f，使 \(h(x) \leqslant p(x)\) 在 E 中成立，且不存在别的线性型 g（E 上的、延拓 f、满足 \(g(x) \leqslant p(x)\) 在 E 中、异于 h 且使 \(g(x) \geqslant h(x)\) 在 S 中成立的）。（考察对 \((V', f')\) 所成之集 F，其中 \(V'\) 是包含 V 的向量子空间，\(f'\) 是 \(V'\) 上延拓 f 的线性型，满足 \(f'(z) \leqslant p(z)\) 在 \(V'\) 中成立，而且不存在别的线性型 \(f''\)（\(V'\) 上的、具有同样性质且使 \(f''(z) \geqslant f'(z)\) 在 \(S \cap V'\) 中成立的）。给 F 排序，并利用 a) 与 S, III, § 2.4 的 th. 2。）
\end{enumerate}

\begin{enumerate}
\def\labelenumi{\arabic{enumi})}
\setcounter{enumi}{6}
\tightlist
\item
  设 \(\mathrm{T}\) 是一个交换幺半群（A, I, § 2.1），带有一个预序关系 \(x \leqslant y\)，使得若 \(x \leqslant y\)，则 \(x + z \leqslant y + z\) 对一切 \(z \in \mathrm{T}\) 成立。一个映射 \(f\)（\(\mathrm{T}\) 到 \(\mathbf{R} \cup \{-\infty\}\) 中的）称为加性的（相应地，次加的、超加的），如果我们有
\end{enumerate}

\[
f (x + y) = f (x) + f (y) \quad (\text { resp. } f (x + y) \leqslant f (x) + f (y), \quad \text { resp. } f (x + y) \geqslant f (x) + f (y))
\]

对任意的 \(x, y\)（\(\mathbf{T}\) 中的）。

\begin{enumerate}
\def\labelenumi{\alph{enumi})}
\item
  若 \(g\) 在 \(\mathbf{T}\) 中是次加的且递增的，则函数 \(h(x) = \inf_{n > 0} g(nx) / n\) 是次加的且递增的；我们有 \(h \leqslant g\)，且 \(h(0) = 0\) 若 \(g(0) \geqslant 0\)。
\item
  在同样的假设下，再设存在两个元素 \(x_{1}, x_{2}\)（\(T\) 的）与两个实数 \(\xi_{1}, \xi_{2}\)，使 \(\xi_{1} < g(x_{1}), \xi_{2} < g(x_{2})\) 且 \(g(x_{1} + x_{2}) < \xi_{1} + \xi_{2}\)。设 \(y_{1}, y_{2}, z_{1}, z_{2}\) 是属于 \(T\) 的四个元素，\(n_{1}, n_{2}\) 是两个整数 \(\geqslant 0\)，并设 \(\alpha_{1}, \alpha_{2}\) 是两个实数，使得
\end{enumerate}

\[
\begin{array}{l} n _ {1} \xi_ {1} + g (z _ {1}) <   \alpha_ {1}, y _ {1} \leqslant n _ {1} x _ {1} + z _ {1} \\ n _ {2} \xi_ {2} + g (z _ {2}) <   \alpha_ {2}, y _ {2} \leqslant n _ {2} x _ {2} + z _ {2}. \end{array}
\]

证明：此时有 \(g(n_2y_1 + n_1y_2) < n_2\alpha_1 + n_1\alpha_2\)。

\begin{enumerate}
\def\labelenumi{\alph{enumi})}
\setcounter{enumi}{2}
\tightlist
\item
  设 \(\omega\) 是 T 上的一个超加函数，满足 \(\omega(0) = 0\)，并设 \(\Omega\) 是 T 上的一个递增次加函数，满足 \(\omega(x) \leqslant \Omega(x)\) 在 T 中成立。证明：存在 T 上的递增加性函数 \(f\)，使 \(\omega(x) \leqslant f(x) \leqslant \Omega(x)\) 在 T 中成立。（注意：递增次加函数 \(g\)（T 上的、使 \(\omega(x) \leqslant g(x) \leqslant \Omega(x)\) 在 T 中成立的）所成之集，对关系 \(\geqslant\) 而言是非空的且归纳的，取该集的一个极小元素作为 \(f\)；利用 \(a\) ) 证明 \(f(0) = 0\)。为证不可能存在 T 的元素对 \((x_1, x_2)\)，使 \(f(x_1 + x_2) < f(x_1) + f(x_2)\)，注意：若 \(\xi_j \in \mathbf{R}\)，且 \(h_j(x) = \inf(n\xi_j + f(y))\)，其中 \(n\) 取遍整数 \(\geqslant 0\) 之集，而 \(y\) 取遍 T 中使 \(x \leqslant nx_j + y\) 的元素所成之集，则 \(h_j\) 在 T 中递增且次加（\(j = 1, 2\)），\(h_j(x_j) \leqslant \xi_j\)，且 \(h_j(x) \leqslant f(x)\) 对一切 \(x \in T\) 成立。然后利用 \(f\) 的定义与 \(b\) ) 得出矛盾。）
\end{enumerate}

¶ * 8) a) 设 K 是一个非离散的赋值除环，其绝对值是超度的，且不是线性紧的（参见 CA, III, § 2, 习题 15）；则存在一个由数 > 0 组成的良序集 I 以及 K 中闭球的一个族 (B(ρ)) \(_{ρ \in I}\)，使得关系 ρ < ρ' 蕴含 B(ρ) ⊂ B(ρ')，B(ρ) 的半径为 ρ，且诸 B(ρ) 的交为空（CA, VI, § 5, 习题 5）。对每个 x ∈ K，存在 ρ ∈ I 使 x ∉ B(ρ)；证明：数 φ(x) = \(|x-y|\)（对某个 \(y \in \mathbf{B}(\rho)\)）既不依赖于 \(y \in \mathbf{B}(\rho)\)，也不依赖于 \(\rho \in \mathbf{I}\)（使 \(x \notin \mathbf{B}(\rho)\) 的）。若 \(\rho \in \mathbf{I}\) 使 \(x \in \mathbf{B}(\rho)\)，则 \(\phi(x) \leqslant \rho\)。在此基础上，对 \((x_1, x_2) \in \mathbf{K}^2\)，令 \(\| (x_1, x_2) \| = |x_1|\) 若 \(x_2 = 0\)，令 \(\| (x_1, x_2) \| = |x_2| \phi(x_2^{-1} x_1)\) 若 \(x_2 \neq 0\)。证明：\(\| (x_1, x_2) \|\) 是 \(\mathbf{K}^2\) 上的一个超范数（I, 第 26 页, 习题 12），并证明不存在 \(\mathbf{K}^2\) 到 \(\mathbf{K} \times \{0\}\) 上的任何范数为 1 的投影。

\begin{enumerate}
\def\labelenumi{\alph{enumi})}
\setcounter{enumi}{1}
\item
  设 K 是一个完备的非离散赋值除环，其绝对值是超度的，且它是线性紧的。设 E 是 K 上带超范数的 2 维向量空间，D 是 E 中一条直线；证明：对一切点 \(x \in \mathbf{E}\)，存在 \(y \in \mathbf{D}\)，使 \(d(x, D) = d(x, y) = \| x - y \|\)（注意 D 与以 \(x\) 为中心的球之交是 D 中的一个球）。
\item
  由 \(a\) ) 与 \(b\) ) 推出：对一个完备非离散赋值除环 \(K\)（其绝对值为超度的），下列性质等价：
\end{enumerate}

\(\alpha)\) K 是线性紧的。

β) 对 K 上每个超赋范向量空间 E、E 的每个向量子空间 F 以及 F 上每个连续线性型 f，存在 E 上的连续线性型 g，它延拓 f 且满足 \(\|g\| = \|f\|\)。（化归到 E 为 2 维的情形并用 b)。）*

